\chapter{Materials and Methods}\label{ch:materials_methods}

% This chapter describes the systematic approach taken to evaluate large language models
% on systems engineering knowledge across multiple question formats (MCQ and OSQ),
% with particular attention to position bias detection and cost-effectiveness analysis.

\section{Research Framework and Overview}\label{sec:research_framework}

This section presents the overall methodological framework that connects the research questions posed in Chapter~\ref{ch:introduction} to the data collection, evaluation, and analysis procedures described in subsequent sections. The research design follows a six-phase pipeline that transforms the base SysEngBench dataset through multiple processing stages, evaluates large language models across different question formats and configurations, and performs comprehensive statistical analysis of the results.

\subsection{Conceptual Model}

The conceptual model for this research is built on three core pillars:

\begin{enumerate}
    \item \textbf{Question Format Diversity}: Evaluating whether different question formats (multiple-choice vs. open-ended short-answer) elicit different levels of understanding from LLMs, analogous to how different assessment types may measure distinct cognitive capabilities in human learners.

    \item \textbf{Robustness Testing}: Systematically varying question presentation (specifically, the position of correct answers in MCQs) to detect systematic biases that may inflate or deflate apparent model performance.

    \item \textbf{Practical Feasibility}: Assessing not just accuracy but also the computational cost and latency implications of different evaluation approaches, particularly for models with extended reasoning capabilities.
\end{enumerate}

These pillars directly address the research questions by providing: (1) comparative performance data across formats, (2) bias-adjusted accuracy metrics, and (3) cost-benefit analysis for practical deployment decisions.

\subsection{Overall Workflow and Data Flow}

Figure~\ref{fig:data_flow_overview} presents the high-level data flow through the six-phase pipeline. Each phase produces well-defined artifacts (datasets, model responses, scores, or analysis outputs) that serve as inputs to subsequent phases, enabling modular execution and independent validation of each processing step.

\begin{figure}[htbp]
    \centering
    % TODO: Include the PlantUML data flow diagram from src/ARCHITECTURE_DIAGRAMS_PLANTUML.md:795-835
    % Convert to PDF using: plantuml -tpdf data_flow.puml
    \includegraphics[width=0.9\textwidth]{figures/data_flow_overview.pdf}
    \caption{Overall data flow through the six-phase evaluation pipeline. Phase 1 downloads and characterizes the base dataset. Phases 2 and 3 run in parallel to generate OSQ conversions and position variants respectively. Phase 4 executes model inference across all configurations. Phase 5 applies LLM-as-a-judge scoring to OSQ responses. Phase 6 performs comprehensive statistical analysis and generates comparative visualizations.}
    \label{fig:data_flow_overview}
\end{figure}

The six phases are:

\begin{description}
    \item[Phase 1: Data Preparation] Downloads the SysEngBench dataset from HuggingFace~\citep{sysengbench_citation} and performs exploratory analysis of question distribution across INCOSE knowledge areas.

    \item[Phase 2: MCQ to OSQ Conversion] Uses a large language model (GPT-4o via OpenRouter) to assess MCQ suitability for open-ended format and performs conversion with rubric generation for questions meeting quality thresholds.

    \item[Phase 3: Position Variant Generation] Creates four variants of the MCQ dataset with correct answers systematically rotated to positions A, B, C, and D to enable position bias detection.

    \item[Phase 4: Model Inference] Executes evaluation of all selected models across six tasks: base MCQ benchmark, four position variants, and OSQ benchmark using lm-evaluation-harness~\citep{lm_eval_harness}.

    \item[Phase 5: LLM-as-a-Judge] Applies multiple judging approaches (binary, rubric-based, multi-dimensional, chain-of-thought) to score OSQ responses using GPT-4o/GPT-4o-mini as judge models~\citep{zheng2023judging}.

    \item[Phase 6: Analysis and Results Processing] Performs statistical analysis including position bias detection, format comparison, tokenomics analysis, and generates publication-ready visualizations.
\end{description}

Phases 2 and 3 are independent and can execute in parallel. All other phases have sequential dependencies as illustrated in Figure~\ref{fig:data_flow_overview}.

\subsection{Experimental Design Principles}

The experimental design adheres to several key principles:

\begin{itemize}
    \item \textbf{Deterministic Evaluation}: All model inference uses temperature=0.0 to enable reproducible results across runs.

    \item \textbf{Semantic Equivalence}: Position variants maintain identical semantic content; only the position of answer choices changes, ensuring that performance differences isolate position bias rather than content variation.

    \item \textbf{Academic Rigor}: Judging methodologies are grounded in peer-reviewed literature (G-Eval~\citep{liu2023geval}, multi-dimensional assessment frameworks~\citep{lin2023multidim}), and statistical testing follows established practices for non-parametric data analysis.

    \item \textbf{Comprehensive Coverage}: Evaluation spans multiple model families (proprietary and open-source), sizes (3B to 405B parameters), and architectural variations (standard vs. reasoning-extended models).

    \item \textbf{Cost Awareness}: All evaluations track token usage and compute cost, enabling economic analysis alongside accuracy assessment.
\end{itemize}

\subsection{Alignment with Research Objectives}

Table~\ref{tab:rq_alignment} maps each research question to the specific phases and analyses that address it.

\begin{table}[htbp]
    \centering
    \caption{Mapping of research questions to methodological phases and analysis components.}
    \label{tab:rq_alignment}
    \begin{tabular}{p{1cm}p{6cm}p{6cm}}
        \toprule
        \textbf{RQ} & \textbf{Research Question} & \textbf{Methodological Components} \\
        \midrule
        RQ1 & How do LLMs perform on systems engineering knowledge assessment? & Phases 1, 4, 6: Base benchmark evaluation, accuracy metrics, cross-category analysis \\
        \addlinespace
        RQ2 & Does question format (MCQ vs. OSQ) affect apparent model capabilities? & Phases 2, 4, 5, 6: OSQ conversion, parallel evaluation, correlation analysis, question-level consistency \\
        \addlinespace
        RQ3 & Are LLMs susceptible to position bias in MCQ evaluations? & Phases 3, 4, 6: Position variants, statistical testing (Chi-Square, Friedman, McNemar), effect size quantification \\
        \addlinespace
        RQ4 & What are the cost-accuracy tradeoffs for different model types? & Phases 4, 6: Token tracking, cost modeling, efficiency metrics, ROI analysis \\
        \bottomrule
    \end{tabular}
\end{table}

%============================================================================
\section{Infrastructure and Technical Environment}\label{sec:infrastructure}

This section details the computational infrastructure, software stack, and technical services used throughout the research pipeline.

\subsection{Computational Infrastructure}

\paragraph{GPU Compute} Model inference (Phase 4) required substantial GPU resources due to the large number of models evaluated and the size of some models (up to 405B parameters). GPU compute was provisioned through RunPod\footnote{\url{https://www.runpod.io/}}, a cloud GPU platform offering on-demand access to NVIDIA A40 and A100 instances.

GPU allocation was determined by model VRAM requirements:
\begin{itemize}
    \item \textbf{Single A40 (48GB VRAM)}: Models up to approximately 43GB (accounting for inference overhead)
    \item \textbf{Dual A40 (96GB VRAM)}: Models up to approximately 90GB
    \item \textbf{Triple A40 (144GB VRAM)}: Largest models requiring up to 141GB
\end{itemize}

Each RunPod instance was configured with:
\begin{itemize}
    \item PyTorch container template (CUDA 12.1, PyTorch 2.1+)
    \item Exposed port 11434 for Ollama API access
    \item Persistent volume storage for model weights (100GB-500GB depending on pod configuration)
    \item SSH access for command-line operations
    \item JupyterLab interface for interactive notebook execution
\end{itemize}

\paragraph{Local LLM Serving} Ollama\footnote{\url{https://ollama.ai/}} (version 0.3.3) served as the local LLM serving layer, providing an OpenAI-compatible API endpoint at \texttt{http://localhost:11434/v1/chat/completions}. Ollama handles model loading, memory management, and efficient batched inference. Models were pulled from the Ollama registry using \texttt{ollama pull <model\_name>} and verified with \texttt{ollama list}.

\paragraph{Development Environment} Primary development occurred on a local workstation running Ubuntu 22.04 LTS with Python 3.11. Jupyter notebooks were used for exploratory analysis and prototyping, while production evaluation scripts were executed on RunPod GPU instances.

\subsection{Software Stack}

The complete software stack is summarized in Table~\ref{tab:software_stack}.

\begin{table}[htbp]
    \centering
    \caption{Software stack components and their roles in the research pipeline.}
    \label{tab:software_stack}
    \begin{tabular}{lll}
        \toprule
        \textbf{Component} & \textbf{Version} & \textbf{Purpose} \\
        \midrule
        Python & 3.11+ & Primary programming language \\
        Pandas & 2.1+ & Data manipulation and CSV processing \\
        NumPy & 1.24+ & Numerical computations \\
        SciPy & 1.11+ & Statistical testing \\
        scikit-learn & 1.3+ & Machine learning utilities \\
        Matplotlib & 3.8+ & Visualization \\
        Seaborn & 0.13+ & Statistical graphics \\
        datasets & 2.14+ & HuggingFace dataset loading \\
        lm-evaluation-harness & 0.4+ & Standardized LLM evaluation \\
        Ollama & 0.3.3 & Local LLM serving \\
        openai & 1.x & API client for judge models \\
        tqdm & 4.66+ & Progress tracking \\
        \bottomrule
    \end{tabular}
\end{table}

\subsection{External Services and APIs}

\paragraph{HuggingFace} The SysEngBench dataset~\citep{sysengbench_citation} was hosted on HuggingFace Hub at \texttt{rabell/SysEngBench} and accessed via the \texttt{datasets} library with authentication token. The converted OSQ dataset was also uploaded to HuggingFace for reproducibility.

\paragraph{OpenRouter} OpenRouter\footnote{\url{https://openrouter.ai/}} provided unified API access to proprietary models (GPT-4o, Claude, Gemini) and larger open-source models not suitable for local execution. This service enabled easy model switching via a single API endpoint while maintaining usage tracking and cost monitoring. OpenRouter was used primarily for:
\begin{itemize}
    \item MCQ to OSQ conversion (Phase 2): GPT-4o
    \item LLM-as-a-judge scoring (Phase 5): GPT-4o and GPT-4o-mini
    \item Supplementary evaluations of proprietary models
\end{itemize}

All API calls used temperature=0.0 for deterministic outputs.

\paragraph{Version Control and Collaboration} All code, notebooks, and analysis scripts were version-controlled using Git and hosted on GitHub at \texttt{ryan-a-bell/dissertation}. LaTeX manuscript development used Overleaf for collaborative editing and automatic compilation.

\subsection{Data Management and Reproducibility}

To ensure reproducibility:
\begin{itemize}
    \item All datasets are versioned and publicly accessible on HuggingFace
    \item Random seeds were fixed where stochastic processes occurred
    \item Model inference used temperature=0.0 for deterministic generation
    \item Complete evaluation logs with timestamps were retained
    \item Environment specifications captured in \texttt{requirements.txt}
    \item All code and notebooks publicly available in the GitHub repository
\end{itemize}

Sensitive credentials (HuggingFace tokens, OpenRouter API keys) were stored in \texttt{.env} files excluded from version control.

%============================================================================
\section{Dataset Preparation}\label{sec:dataset_construction}

Dataset preparation consists of three parallel workflows: acquiring and characterizing the base SysEngBench dataset (Phase 1), converting suitable MCQs to open-ended format with grading rubrics (Phase 2), and generating position-rotated variants for bias testing (Phase 3).

\subsection{Base Dataset Acquisition (Phase 1)}\label{sec:phase1}

\subsubsection{SysEngBench Download and Schema}

SysEngBench~\citep{sysengbench_citation} is a multiple-choice question benchmark for systems engineering knowledge, aligned to the INCOSE Systems Engineering Handbook. The dataset was downloaded from HuggingFace (\texttt{rabell/SysEngBench}) using authenticated API access. The test split contains 1,144 questions spanning six major INCOSE knowledge areas and 20 sub-categories.

Each question record includes the following fields:
\begin{itemize}
    \item \texttt{question\_id}: Unique identifier
    \item \texttt{tags}: Comma-separated metadata tags
    \item \texttt{incose\_category}: INCOSE Handbook categorization (major category / sub-category)
    \item \texttt{question}: Question text
    \item \texttt{choiceA}, \texttt{choiceB}, \texttt{choiceC}, \texttt{choiceD}: Four answer options
    \item \texttt{answer}: Correct answer letter (A/B/C/D)
    \item \texttt{label}: Numeric label (0/1/2/3 corresponding to A/B/C/D)
    \item \texttt{justification}: Explanation of correct answer
\end{itemize}

The dataset was saved locally as \texttt{sysengbench.csv} for subsequent processing.

Figure~\ref{fig:phase1_download} shows the download workflow.

\begin{figure}[htbp]
    \centering
    % TODO: Include PlantUML diagram from src/ARCHITECTURE_DIAGRAMS_PLANTUML.md:13-36
    \includegraphics[width=0.6\textwidth]{figures/phase1_download.pdf}
    \caption{Phase 1 data download workflow. The SysEngBench dataset is loaded from HuggingFace with authentication, the test split extracted, converted to Pandas DataFrame format, and saved locally.}
    \label{fig:phase1_download}
\end{figure}

\subsubsection{Dataset Characterization}

Exploratory analysis examined the distribution of questions across INCOSE categories and sub-categories. Questions were categorized into six major knowledge areas as shown in Table~\ref{tab:incose_distribution}.

\begin{table}[htbp]
    \centering
    \caption{Distribution of SysEngBench questions across INCOSE Handbook categories.}
    \label{tab:incose_distribution}
    \begin{tabular}{lrr}
        \toprule
        \textbf{INCOSE Category} & \textbf{Questions} & \textbf{Percentage} \\
        \midrule
        Cross-Cutting Systems Engineering Methods & 276 & 24.1\% \\
        Specialty Engineering Activities & 419 & 36.6\% \\
        Systems Engineering Management & XXX & XX.X\% \\
        Technical Processes & XXX & XX.X\% \\
        % TODO: Fill in remaining categories from actual data
        \bottomrule
        \textbf{Total} & \textbf{1,144} & \textbf{100.0\%} \\
    \end{tabular}
\end{table}

The 20 sub-categories were also analyzed, with distributions visualized using bar charts, pie charts, and stacked visualizations (see Appendix~\ref{app:dataset_viz}). This characterization ensures that subsequent analyses can account for category-level performance variations and that the dataset provides reasonable coverage across the systems engineering knowledge domain.











%%%%%%%%%%%%%% STOPPED HERE... %%%% 

%%%%%%%%%%%%%% PICK BACK UP HERE FOR OVERLEAD INTEGRATION %%%%%%%%%%%%%%










\subsection{MCQ to OSQ Conversion (Phase 2)}\label{sec:phase2}

A significant methodological contribution of this research is the systematic conversion of suitable multiple-choice questions to open-ended short-answer (OSQ) format using a two-stage LLM-based pipeline. This conversion enables comparative analysis of model performance across question formats while maintaining alignment to the same underlying knowledge construct.

\subsubsection{Two-Stage Conversion Pipeline}

The conversion process consists of two sequential stages: (1) suitability classification, and (2) conversion with rubric generation. Figure~\ref{fig:phase2_conversion} illustrates the complete pipeline.

\begin{figure}[htbp]
    \centering
    % TODO: Include PlantUML diagram from src/ARCHITECTURE_DIAGRAMS_PLANTUML.md:108-173
    \includegraphics[width=0.9\textwidth]{figures/phase2_conversion.pdf}
    \caption{Phase 2 two-stage MCQ to OSQ conversion pipeline. Stage 1 classifies each MCQ for OSQ suitability across five dimensions. Questions scoring $\geq 7/10$ proceed to Stage 2, where conversion generates OSQ prompt, expected answer, grading rubric, and Bloom's taxonomy classification.}
    \label{fig:phase2_conversion}
\end{figure}

\paragraph{Stage 1: Suitability Classification}

Not all multiple-choice questions are suitable for conversion to open-ended format. Questions requiring visual diagrams, those with highly ambiguous correct answers, or those where the distractor options provide essential context may not translate well.

Each of the 1,144 MCQs was evaluated by GPT-4o (via OpenRouter, temperature=0.0) across five dimensions:

\begin{enumerate}
    \item \textbf{Assessability}: Can the answer be evaluated objectively in open-ended format?
    \item \textbf{Canonical Answer Clarity}: Is there a well-defined expected answer?
    \item \textbf{Ambiguity}: Is the question clear without the MCQ structure?
    \item \textbf{Scope Appropriateness}: Is the expected answer reasonable in length for short-answer format?
    \item \textbf{Intent Preservation}: Can the original knowledge construct be maintained without MCQ options?
\end{enumerate}

The LLM returned a structured JSON response containing:
\begin{verbatim}
{
  "suitability_score": <1-10>,
  "justification": "<reasoning>",
  "dimension_scores": {
    "assessability": <1-10>,
    "canonical_answer": <1-10>,
    "ambiguity": <1-10>,
    "scope": <1-10>,
    "intent_preservation": <1-10>
  }
}
\end{verbatim}

Questions with overall suitability score $\geq 7$ were marked for conversion. Of 1,144 questions, 845 (73.9\%) met this threshold.

The complete classification prompt is provided in Appendix~\ref{app:conversion_prompts}.

\paragraph{Stage 2: Conversion with Rubric Generation}

For each suitable MCQ, GPT-4o performed conversion while preserving the original knowledge construct, difficulty level, and cognitive demand. The conversion process:

\begin{itemize}
    \item Removed MCQ-specific artifacts (e.g., ``which of the following'', explicit A/B/C/D options)
    \item Maintained the original intent (e.g., definitional recall remains recall, not application)
    \item Preserved any numerical calculations or technical specifications
    \item Generated a canonical expected answer based on the original correct choice
    \item Created a detailed grading rubric with three levels:
    \begin{itemize}
        \item \textbf{Full credit}: Criteria for complete correct answer
        \item \textbf{Partial credit}: Criteria for partially correct answer
        \item \textbf{No credit}: Conditions warranting no credit
    \end{itemize}
    \item Assigned Bloom's taxonomy level (Remember, Understand, Apply, Analyze, Evaluate, Create)
\end{itemize}

The conversion prompt explicitly instructed the LLM to avoid ``improving'' the question or making it more challenging, maintaining parity with the original MCQ construct.

The LLM returned structured JSON:
\begin{verbatim}
{
  "osq_prompt": "<open-ended question>",
  "expected_answer": "<canonical answer>",
  "rubric": {
    "full_credit": "<criteria>",
    "partial_credit": "<criteria>",
    "no_credit": "<criteria>"
  },
  "blooms_level": "<taxonomy level>",
  "blooms_justification": "<reasoning>"
}
\end{verbatim}

Example conversion in Table~\ref{tab:conversion_example}.

\begin{table}[htbp]
    \centering
    \caption{Example MCQ to OSQ conversion showing original question, converted format, and generated rubric.}
    \label{tab:conversion_example}
    \small
    \begin{tabular}{p{3cm}p{10cm}}
        \toprule
        \textbf{Field} & \textbf{Content} \\
        \midrule
        \textbf{Original MCQ} &
        \begin{minipage}[t]{10cm}
        What is the primary purpose of stakeholder analysis in systems engineering?
        \begin{itemize}[noitemsep,topsep=0pt]
            \item[A.] To identify project funding sources
            \item[B.] To understand stakeholder needs and concerns
            \item[C.] To assign project responsibilities
            \item[D.] To develop marketing strategies
        \end{itemize}
        \textbf{Answer}: B
        \end{minipage} \\
        \addlinespace
        \textbf{OSQ Prompt} & What is the primary purpose of stakeholder analysis in systems engineering? \\
        \addlinespace
        \textbf{Expected Answer} & To understand stakeholder needs and concerns \\
        \addlinespace
        \textbf{Full Credit} & Response identifies understanding stakeholder needs, concerns, expectations, or requirements as primary purpose \\
        \addlinespace
        \textbf{Partial Credit} & Response mentions stakeholders but focuses on secondary purposes (e.g., communication planning) \\
        \addlinespace
        \textbf{No Credit} & Response focuses on unrelated activities (funding, marketing) or misidentifies purpose \\
        \addlinespace
        \textbf{Bloom's Level} & Remember (definitional recall) \\
        \bottomrule
    \end{tabular}
\end{table}

\subsubsection{Quality Control and Validation}

Several quality control mechanisms ensured conversion fidelity:

\begin{itemize}
    \item \textbf{Deterministic Generation}: Temperature=0.0 ensured consistent outputs across runs
    \item \textbf{Comprehensive Logging}: All prompts and responses were logged with timestamps to separate artifact directories
    \item \textbf{JSON Validation}: Retry logic handled parsing failures
    \item \textbf{Manual Spot-Checks}: A random sample of 50 conversions was manually reviewed to verify semantic preservation and rubric quality
\end{itemize}

\subsubsection{Output Datasets}

Phase 2 produced three datasets:

\begin{enumerate}
    \item \textbf{sysengbench\_osq.csv}: All 1,144 questions with additional OSQ columns (osq\_prompt, expected\_answer, rubric components, Bloom's level, suitability score)
    \item \textbf{sysengbench\_osq\_filtered.csv}: Only the 845 suitable questions (osq\_suitable == ``Yes'')
    \item \textbf{sysengbench\_osq.jsonl}: JSONL format for HuggingFace upload and lm-evaluation-harness compatibility
\end{enumerate}

All three datasets were version-controlled and uploaded to HuggingFace for reproducibility.

\subsection{Position Variant Generation (Phase 3)}\label{sec:phase3}

To detect position bias in model responses to multiple-choice questions, four variants of the SysEngBench dataset were generated with correct answers systematically rotated to positions A, B, C, and D. This approach enables statistical testing for systematic preference for particular answer positions independent of question content.

\subsubsection{Circular Rotation Algorithm}

The rotation algorithm maintains semantic equivalence by circularly shifting the choice list without altering content. Given a question with current correct answer at position $p_{\text{current}}$ and target position $p_{\text{target}}$, the shift amount $k$ is calculated as:

\begin{equation}
k = (p_{\text{target}} - p_{\text{current}}) \bmod 4
\end{equation}

The choice list is then rotated:
\begin{equation}
\text{choices}_{\text{rotated}} = \text{choices}[-k:] + \text{choices}[:-k]
\end{equation}

For example, to move the correct answer from position B (index 1) to position A (index 0):
\begin{align*}
k &= (0 - 1) \bmod 4 = 3 \\
[A, \mathbf{B}, C, D] &\rightarrow [B, C, D, A]
\end{align*}

Figure~\ref{fig:phase3_rotation} illustrates the rotation process.

\begin{figure}[htbp]
    \centering
    % TODO: Include PlantUML diagram from src/ARCHITECTURE_DIAGRAMS_PLANTUML.md:208-257
    \includegraphics[width=0.9\textwidth]{figures/phase3_rotation.pdf}
    \caption{Phase 3 position variant generation using circular rotation. Four parallel workflows create variants with correct answers at positions A, B, C, and D respectively. Each variant contains all 1,144 questions with semantically equivalent but position-rotated choices.}
    \label{fig:phase3_rotation}
\end{figure}

\subsubsection{Variant Characteristics}

Four variants were generated:

\begin{itemize}
    \item \textbf{sysengbench\_a.csv}: Correct answer always at position A (label=0)
    \item \textbf{sysengbench\_b.csv}: Correct answer always at position B (label=1)
    \item \textbf{sysengbench\_c.csv}: Correct answer always at position C (label=2)
    \item \textbf{sysengbench\_d.csv}: Correct answer always at position D (label=3)
\end{itemize}

Each variant contains all 1,144 questions. The schema is identical to the base dataset. Semantic content is preserved; only choice ordering differs.

\subsubsection{Verification}

Verification ensured correct rotation by:
\begin{enumerate}
    \item Previewing the first row of each variant to confirm correct answer position
    \item Verifying that choice content (after rotation) matches expected pattern
    \item Confirming that all 1,144 questions appear in each variant
\end{enumerate}

These variants enable Phase 6 position bias analysis (Section~\ref{sec:position_bias_analysis}) by comparing model accuracy across variants. If no position bias exists, accuracy should be uniform across all four variants. Systematic differences indicate position preference.

%============================================================================
\section{Evaluation Procedures}\label{sec:evaluation_methods_computations}

Model evaluation occurs in Phase 4 (inference) and Phase 5 (LLM-as-a-judge scoring). This section describes the evaluation configurations, prompt formatting, and judging methodologies.

\subsection{MCQ Evaluation Configuration (Phase 4)}\label{sec:mcq_eval}

\subsubsection{Model Selection}

Models were selected to represent diverse architectural families, parameter scales, and licensing types:

\begin{itemize}
    \item \textbf{Open-source models} (via Ollama): Llama 3.2 (3B, 7B, 70B), Qwen 2.5 (7B, 14B, 32B, 72B), Mistral (7B), DeepSeek-R1 (7B, reasoning-extended), and others
    \item \textbf{Proprietary models} (via OpenRouter): GPT-4o, GPT-4o-mini, Claude 3.5 Sonnet, Gemini 1.5 Pro
    \item \textbf{Reasoning-extended models}: DeepSeek-R1, QwQ-32B-Preview
\end{itemize}

Models were filtered by VRAM availability on RunPod instances. A complete model list is provided in Appendix~\ref{app:model_list}.

\subsubsection{Evaluation Framework}

All MCQ evaluations used lm-evaluation-harness~\citep{lm_eval_harness}, a standardized framework for LLM benchmarking. This framework provides:
\begin{itemize}
    \item Consistent prompt formatting across models
    \item Automatic metric calculation (accuracy, per-category accuracy)
    \item Structured output (JSON results, JSONL sample logs)
    \item Support for custom tasks via YAML configuration
\end{itemize}

\subsubsection{Task Configuration}

Five MCQ tasks were defined:
\begin{enumerate}
    \item \texttt{sysengbench}: Base benchmark (original answer distribution)
    \item \texttt{sysengbench-a}: Variant with answers at position A
    \item \texttt{sysengbench-b}: Variant with answers at position B
    \item \texttt{sysengbench-c}: Variant with answers at position C
    \item \texttt{sysengbench-d}: Variant with answers at position D
\end{enumerate}

Each task was configured with the YAML specification shown in Listing~\ref{lst:mcq_config}.

\begin{lstlisting}[caption={lm-evaluation-harness task configuration for MCQ evaluation.},label={lst:mcq_config},language=yaml,basicstyle=\small\ttfamily]
task: sysengbench
dataset_path: rabell/SysEngBench
dataset_name: default
output_type: generate_until
doc_to_text: |
  Given the following question and four candidate answers (A, B, C, D),
  choose the best answer.

  Question: {{question}}
  A. {{choiceA}}
  B. {{choiceB}}
  C. {{choiceC}}
  D. {{choiceD}}

  Answer:
doc_to_target: "{{answer}}"
generation_kwargs:
  temperature: 0.0
  max_tokens: 20
  until: ["</s>", "\n"]
filter_list:
  - name: regex
    regex_pattern: "([ABCD])"
  - name: take_first
metric_list:
  - metric: exact_match
    aggregation: mean
    higher_is_better: true
\end{lstlisting}

Key configuration choices:
\begin{itemize}
    \item \textbf{Temperature=0.0}: Ensures deterministic, reproducible generation
    \item \textbf{Max tokens=20}: Sufficient for single-letter answer plus brief explanation
    \item \textbf{Stop sequences}: \texttt{</s>} (end-of-sequence token) and \texttt{\textbackslash n} (newline)
    \item \textbf{Regex filter}: Extracts first occurrence of A, B, C, or D from response
    \item \textbf{Exact match metric}: Binary correct/incorrect at question level
\end{itemize}

\subsubsection{Batch Inference Process}

Model inference followed the workflow in Figure~\ref{fig:phase4_inference}.

\begin{figure}[htbp]
    \centering
    % TODO: Include PlantUML diagram from src/ARCHITECTURE_DIAGRAMS_PLANTUML.md:365-432
    \includegraphics[width=0.8\textwidth]{figures/phase4_inference.pdf}
    \caption{Phase 4 batch inference workflow with error handling. The process iterates through all models, checks server health, executes lm-evaluation-harness, and logs success/failure with automatic retry on server errors.}
    \label{fig:phase4_inference}
\end{figure}

For each model in the model list:
\begin{enumerate}
    \item Check if Ollama server is running; restart if necessary (max 3 retries)
    \item Build lm\_eval command:
    \begin{verbatim}
lm_eval --model local-chat-completions \
        --model_args base_url=http://localhost:11434/v1,model=<model_name> \
        --tasks <task_name> \
        --output_path output/<task_name>/ \
        --log_samples \
        --batch_size auto
    \end{verbatim}
    \item Execute via subprocess with output logging
    \item On success: log completion timestamp
    \item On failure: check if server error; if yes, restart server and retry once
    \item Append failed models to \texttt{failed\_models} list
    \item Write progress to timestamped log file
\end{enumerate}

This process was repeated for all five MCQ tasks, generating 5 evaluation runs per model.

\subsubsection{Output Structure}

For each task and model combination, lm-evaluation-harness generated:
\begin{itemize}
    \item \texttt{results\_<timestamp>.json}: Aggregate metrics (overall accuracy, per-category accuracy)
    \item \texttt{samples\_<timestamp>.jsonl}: Per-question details (doc\_id, question, model response, target, correctness)
\end{itemize}

Complete output structure:
\begin{verbatim}
output/
├── sysengbench/
│   ├── <model1>/
│   │   ├── results_<timestamp>.json
│   │   └── samples_<timestamp>.jsonl
│   ├── <model2>/
│   └── ...
├── sysengbench-a/
├── sysengbench-b/
├── sysengbench-c/
└── sysengbench-d/
\end{verbatim}

\subsection{OSQ Evaluation Configuration (Phase 4)}\label{sec:osq_eval}

\subsubsection{Task Configuration}

OSQ evaluation used the \texttt{sysengbench-osq} task with configuration shown in Listing~\ref{lst:osq_config}.

\begin{lstlisting}[caption={lm-evaluation-harness task configuration for OSQ evaluation.},label={lst:osq_config},language=yaml,basicstyle=\small\ttfamily]
task: sysengbench-osq
dataset_path: rabell/SysEngBench-OSQ
dataset_name: default
output_type: generate_until
doc_to_text: "{{osq_prompt}}"
generation_kwargs:
  temperature: 0.0
  max_tokens: 2000  # Varies by model context window
  until: ["</s>"]
# No automatic metric calculation - responses logged for Phase 5 judging
\end{lstlisting}

Key differences from MCQ configuration:
\begin{itemize}
    \item \textbf{Higher max\_tokens}: 2000-10000 depending on model, allowing for detailed explanations
    \item \textbf{No filtering}: Full response text retained
    \item \textbf{No automatic metrics}: Accuracy cannot be computed via string matching; requires semantic judging (Phase 5)
\end{itemize}

\subsubsection{Output for Judging}

OSQ evaluation generated \texttt{samples\_<timestamp>.jsonl} files containing:
\begin{verbatim}
{
  "doc_id": 42,
  "doc": {
    "question_id": "...",
    "osq_prompt": "...",
    "expected_answer": "...",
    "rubric_full_credit": "...",
    "rubric_partial_credit": "...",
    "rubric_no_credit": "..."
  },
  "target": null,
  "model_response": "<model's open-ended answer>",
  "filtered_response": "<model's open-ended answer>"
}
\end{verbatim}

These files served as input to Phase 5 LLM-as-a-judge scoring.


%%%%%%%%%%%%%%%% STOPPPED HERE %%%%%%%%%%%%%%%% PICK BACK UP HERE %%%%%%%%%%%%%%



\subsection{LLM-as-a-Judge Scoring (Phase 5)}\label{sec:llm_judge}

Scoring open-ended responses requires semantic understanding rather than simple string matching. This research employs LLM-as-a-judge~\citep{zheng2023judging}, a methodology where a capable language model evaluates responses based on rubrics and expected answers.

\subsubsection{Judge Model Selection}

Two judge models were used:
\begin{itemize}
    \item \textbf{GPT-4o} (primary judge): High-capability model with strong instruction following
    \item \textbf{GPT-4o-mini} (secondary judge): Lower-cost alternative for cost-sensitivity analysis
\end{itemize}

Both accessed via OpenRouter with temperature=0.0.

\subsubsection{Four Judging Approaches}

To validate scoring robustness, four judging approaches were implemented based on peer-reviewed methodologies:

\paragraph{1. Binary Correct/Incorrect}
Simplest baseline: judge returns ``correct'' or ``incorrect'' with brief justification.

\paragraph{2. Rubric-Based Judging}
Uses the grading rubric generated in Phase 2 (Section~\ref{sec:phase2}). Judge assigns:
\begin{itemize}
    \item \textbf{Full credit} (1.0): Response meets full credit criteria
    \item \textbf{Partial credit} (0.5): Response meets partial credit criteria
    \item \textbf{No credit} (0.0): Response meets no credit criteria
\end{itemize}

Prompt structure (abridged; full prompt in Appendix~\ref{app:judge_prompts}):
\begin{verbatim}
You are grading a systems engineering question response.

Question: {osq_prompt}
Expected Answer: {expected_answer}

Rubric:
- Full credit: {rubric_full_credit}
- Partial credit: {rubric_partial_credit}
- No credit: {rubric_no_credit}

Student Response: {model_response}

Assign a score (0.0, 0.5, or 1.0) and provide reasoning.
\end{verbatim}

\paragraph{3. Multi-Dimensional Scoring}
Based on Lin \& Chen (2023)~\citep{lin2023multidim}, this approach evaluates five dimensions on 0-10 scale:

\begin{enumerate}
    \item \textbf{Technical Accuracy}: Correctness of facts, principles, terminology
    \item \textbf{Conceptual Understanding}: Depth of systems thinking demonstrated
    \item \textbf{Completeness}: Coverage of required answer elements
    \item \textbf{Clarity}: Communication quality and organization
    \item \textbf{Professional Relevance}: Real-world applicability
\end{enumerate}

Aggregate score computed as:
\begin{equation}
\text{Score}_{\text{multi}} = \frac{1}{5} \sum_{i=1}^{5} \frac{\text{Dimension}_i}{10}
\end{equation}

This produces a normalized score in [0, 1].

\paragraph{4. Chain-of-Thought (CoT) Scoring}
Based on G-Eval~\citep{liu2023geval} and Zheng et al.~\citep{zheng2023judging}, the judge is prompted to:
\begin{enumerate}
    \item Reason step-by-step about response quality
    \item Identify specific strengths and weaknesses
    \item Provide final score with justification
\end{enumerate}

This approach shows higher human agreement rates (85-98\%) in prior work~\citep{zheng2023judging}.

\subsubsection{Judging Process}

Figure~\ref{fig:phase5_judge} illustrates the judging workflow.

\begin{figure}[htbp]
    \centering
    % TODO: Include PlantUML diagram from src/ARCHITECTURE_DIAGRAMS_PLANTUML.md:493-556
    \includegraphics[width=0.8\textwidth]{figures/phase5_judge.pdf}
    \caption{Phase 5 LLM-as-a-judge workflow. The process loads OSQ responses, selects judging approach, iterates through samples with resume capability, and writes judgments incrementally to JSONL.}
    \label{fig:phase5_judge}
\end{figure}

For each model's OSQ responses:
\begin{enumerate}
    \item Load samples JSONL from Phase 4
    \item Check for existing judged output; if present, resume from last processed sample
    \item For each unprocessed sample:
    \begin{enumerate}
        \item Build judge prompt with question, expected answer, rubric/criteria, and model response
        \item Call judge model API (temperature=0.0)
        \item Parse JSON response containing score, reasoning, and dimension breakdown (if applicable)
        \item Append to \texttt{samples\_judged.jsonl} immediately (crash-safe incremental writing)
        \item Update progress bar
    \end{enumerate}
    \item Save final organized output matching Phase 4 directory structure
\end{enumerate}

\subsubsection{Output Format}

Judged samples are stored in JSONL format:
\begin{verbatim}
{
  "task": "sysengbench-osq",
  "doc_id": 42,
  "question": "...",
  "expected_answer": "...",
  "model_response": "...",
  "judge_model": "gpt-4o",
  "judging_approach": "multi_dimensional",
  "judgment": {
    "score": 0.82,
    "reasoning": "...",
    "technical_accuracy": 9,
    "conceptual_understanding": 8,
    "completeness": 9,
    "clarity": 7,
    "professional_relevance": 8
  }
}
\end{verbatim}

\subsubsection{Reliability and Bias Mitigation}

Several strategies mitigate known biases in LLM judging:
\begin{itemize}
    \item \textbf{Temperature=0.0}: Reduces stochasticity
    \item \textbf{Explicit rubrics}: Grounds evaluation in concrete criteria
    \item \textbf{Position bias mitigation}: Expected answer and model response order randomized (when applicable)
    \item \textbf{Length bias mitigation}: Rubrics emphasize correctness over verbosity
    \item \textbf{Multiple approaches}: Four judging methods enable cross-validation
\end{itemize}

Inter-judge reliability (between GPT-4o and GPT-4o-mini) is analyzed in Phase 6 (Section~\ref{sec:judge_agreement}).

%============================================================================
\section{Analysis Metrics and Methods (Phase 6)}\label{sec:analysis_metrics_methods}

Phase 6 aggregates results from Phases 4 and 5, performing comprehensive statistical analysis across three primary dimensions: position bias, format comparison (MCQ vs. OSQ), and tokenomics/cost-effectiveness. This section describes the metrics, statistical tests, and visualization approaches for each analysis dimension.

\subsection{Position Bias Analysis}\label{sec:position_bias_analysis}

Position bias refers to systematic preference for particular answer positions (A, B, C, D) independent of question content. Such bias can inflate or deflate apparent model accuracy depending on the distribution of correct answers in the benchmark.

\subsubsection{Metrics}

For each model, accuracy is calculated separately for each position variant:
\begin{align}
\text{Acc}_A &= \frac{\text{\# correct on variant A}}{1144} \\
\text{Acc}_B &= \frac{\text{\# correct on variant B}}{1144} \\
\text{Acc}_C &= \frac{\text{\# correct on variant C}}{1144} \\
\text{Acc}_D &= \frac{\text{\# correct on variant D}}{1144}
\end{align}

Under the null hypothesis of no position bias, $\text{Acc}_A = \text{Acc}_B = \text{Acc}_C = \text{Acc}_D$.

Deviation from expected uniform distribution is quantified as:
\begin{equation}
\Delta_p = \text{Acc}_p - \frac{1}{4}\sum_{i \in \{A,B,C,D\}} \text{Acc}_i
\end{equation}

\subsubsection{Statistical Tests}

Multiple statistical tests are applied to assess position bias significance:

\paragraph{Chi-Square Test for Independence}
Tests whether response distribution across positions differs from expected:
\begin{equation}
\chi^2 = \sum_{p \in \{A,B,C,D\}} \frac{(O_p - E_p)^2}{E_p}
\end{equation}
where $O_p$ is observed correct responses at position $p$ and $E_p$ is expected under uniform distribution.

Degrees of freedom: $df = 3$. Significance threshold: $\alpha = 0.05$.

\paragraph{Kruskal-Wallis H Test}
Non-parametric alternative to one-way ANOVA, testing whether accuracy distributions differ across positions without assuming normality:
\begin{equation}
H = \frac{12}{N(N+1)} \sum_{p=1}^{4} \frac{R_p^2}{n_p} - 3(N+1)
\end{equation}
where $R_p$ is sum of ranks for position $p$, $n_p$ is sample size (1144), and $N = 4 \times 1144$.

\paragraph{Friedman Test}
Non-parametric repeated measures test treating each question as a ``subject'' evaluated across four ``conditions'' (positions):
\begin{equation}
\chi_F^2 = \frac{12}{nk(k+1)} \sum_{p=1}^{4} R_p^2 - 3n(k+1)
\end{equation}
where $n = 1144$ questions, $k = 4$ positions, $R_p$ is sum of ranks for position $p$.

\paragraph{Pairwise McNemar Tests}
For each pair of positions $(p_1, p_2)$, McNemar's test assesses whether the model's responses differ systematically:
\begin{equation}
\chi^2_{\text{McNemar}} = \frac{(b - c)^2}{b + c}
\end{equation}
where $b$ is count of questions correct on $p_1$ but incorrect on $p_2$, and $c$ is the reverse.

Bonferroni correction applied for multiple comparisons: $\alpha_{\text{adj}} = 0.05 / 6 = 0.0083$ for six pairwise tests.

\paragraph{Cramér's V (Effect Size)}
Quantifies effect size for Chi-Square test:
\begin{equation}
V = \sqrt{\frac{\chi^2}{N \cdot \min(r-1, c-1)}}
\end{equation}

Interpretation: $V < 0.1$ (negligible), $0.1 \leq V < 0.3$ (small), $0.3 \leq V < 0.5$ (medium), $V \geq 0.5$ (large)~\citep{cohen1988}.

\subsubsection{Normality Testing}

Before parametric tests, normality of accuracy distributions is assessed using:
\begin{itemize}
    \item \textbf{Shapiro-Wilk test}: $W$ statistic, $p < 0.05$ indicates non-normality
    \item \textbf{Jarque-Bera test}: Based on skewness and kurtosis
    \item \textbf{Q-Q plots}: Visual assessment of deviation from normal distribution
\end{itemize}

If normality is violated, non-parametric tests (Kruskal-Wallis, Friedman) are preferred.

\subsubsection{Visualizations}

Position bias is visualized via:
\begin{itemize}
    \item Bar charts: Accuracy by position for each model
    \item Heatmaps: Position preference across all models
    \item Deviation plots: $\Delta_p$ for each model and position
    \item Statistical significance annotations: Asterisks indicating $p < 0.05$, $p < 0.01$, $p < 0.001$
\end{itemize}

Example visualization shown in Figure~\ref{fig:position_bias_viz} (placeholder; actual figure generated in Phase 6).

\begin{figure}[htbp]
    \centering
    % TODO: Include actual Phase 6 position bias visualization
    \includegraphics[width=0.9\textwidth]{figures/position_bias_example.pdf}
    \caption{Position bias heatmap showing accuracy by answer position across evaluated models. Darker colors indicate higher accuracy. Statistical significance markers denote positions with significant deviation from mean accuracy (Chi-Square test, $p < 0.05$).}
    \label{fig:position_bias_viz}
\end{figure}

\subsection{Format Comparison: MCQ vs. OSQ}\label{sec:format_comparison}

This analysis assesses whether question format (multiple-choice vs. open-ended) affects apparent model capabilities, addressing RQ2.

\subsubsection{Metrics}

For each model:
\begin{align}
\text{Acc}_{\text{MCQ}} &= \frac{\text{\# correct MCQ}}{1144} \\
\text{Acc}_{\text{OSQ}} &= \frac{\sum \text{OSQ scores}}{845}
\end{align}

Note: OSQ accuracy computed over 845 suitable questions (from Phase 2), while MCQ uses full 1144. For fair comparison, MCQ accuracy is also computed on the subset of 845 questions.

\paragraph{Overall Comparison}
\begin{itemize}
    \item Mean accuracy difference: $\overline{\text{Acc}}_{\text{MCQ}} - \overline{\text{Acc}}_{\text{OSQ}}$
    \item Paired t-test (if normality holds) or Wilcoxon signed-rank test (non-parametric)
\end{itemize}

\paragraph{Correlation Analysis}
Pearson correlation coefficient between MCQ and OSQ accuracy across models:
\begin{equation}
r = \frac{\sum (x_i - \bar{x})(y_i - \bar{y})}{\sqrt{\sum (x_i - \bar{x})^2 \sum (y_i - \bar{y})^2}}
\end{equation}
where $x_i = \text{Acc}_{\text{MCQ}}^{(i)}$ and $y_i = \text{Acc}_{\text{OSQ}}^{(i)}$ for model $i$.

Strong positive correlation ($r > 0.8$) suggests formats measure similar construct; weak correlation ($r < 0.5$) suggests different capabilities.

\paragraph{Question-Level Consistency}
For each question $q$ in the 845 OSQ set, define:
\begin{equation}
\text{Consistency}_q = \frac{\# \text{models correct on both formats}}{\# \text{models evaluated}}
\end{equation}

High consistency ($> 0.8$) indicates question is format-independent; low consistency ($< 0.5$) indicates format-sensitive.

\subsubsection{Statistical Tests}

\paragraph{Paired t-test}
Null hypothesis: $\mu_{\text{MCQ}} = \mu_{\text{OSQ}}$

Test statistic:
\begin{equation}
t = \frac{\bar{d}}{s_d / \sqrt{n}}
\end{equation}
where $\bar{d}$ is mean difference, $s_d$ is standard deviation of differences, $n$ is number of models.

If normality violated (Shapiro-Wilk $p < 0.05$), use Wilcoxon signed-rank test instead.

\paragraph{Effect Size (Cohen's d)}
\begin{equation}
d = \frac{\bar{x}_{\text{MCQ}} - \bar{x}_{\text{OSQ}}}{s_{\text{pooled}}}
\end{equation}
where
\begin{equation}
s_{\text{pooled}} = \sqrt{\frac{(n_1 - 1)s_1^2 + (n_2 - 1)s_2^2}{n_1 + n_2 - 2}}
\end{equation}

Interpretation: $|d| < 0.2$ (negligible), $0.2 \leq |d| < 0.5$ (small), $0.5 \leq |d| < 0.8$ (medium), $|d| \geq 0.8$ (large)~\citep{cohen1988}.

\paragraph{Bootstrap Confidence Intervals}
Non-parametric bootstrap (10,000 resamples) estimates confidence intervals for mean difference:
\begin{enumerate}
    \item Resample models with replacement
    \item Calculate $\bar{d}^*$ for resample
    \item Repeat 10,000 times
    \item Compute 95\% CI as 2.5th and 97.5th percentiles
\end{enumerate}

\subsubsection{Visualizations}

\begin{itemize}
    \item Scatter plots: MCQ accuracy (x-axis) vs. OSQ accuracy (y-axis), one point per model
    \item Correlation line with $r$ value annotated
    \item Side-by-side bar charts: MCQ and OSQ accuracy for each model
    \item Bland-Altman plot: Mean accuracy vs. difference (MCQ - OSQ)
    \item Question-level consistency histogram
\end{itemize}

Example scatter plot in Figure~\ref{fig:mcq_osq_scatter}.

\begin{figure}[htbp]
    \centering
    % TODO: Include actual Phase 6 MCQ vs OSQ scatter plot
    \includegraphics[width=0.7\textwidth]{figures/mcq_osq_scatter.pdf}
    \caption{Scatter plot of MCQ accuracy vs. OSQ accuracy across evaluated models. Dashed line represents $y = x$ (perfect agreement). Pearson correlation coefficient $r$ is annotated. Points above the line indicate models performing better on MCQ; points below indicate better OSQ performance.}
    \label{fig:mcq_osq_scatter}
\end{figure}

\subsection{Tokenomics and Cost-Effectiveness Analysis}\label{sec:tokenomics}

This analysis quantifies the cost-accuracy tradeoff, particularly for reasoning-extended models that generate substantially more tokens but may achieve higher accuracy (addressing RQ4).

\subsubsection{Metrics}

For each model $m$:

\paragraph{Token Usage}
\begin{align}
\text{Tokens}_{\text{avg}}^{(m)} &= \frac{1}{N} \sum_{i=1}^{N} \text{tokens}(r_i^{(m)}) \\
\text{Tokens}_{\text{total}}^{(m)} &= \sum_{i=1}^{N} \text{tokens}(r_i^{(m)})
\end{align}
where $r_i^{(m)}$ is response to question $i$ by model $m$, and $N = 1144$ for MCQ or $N = 845$ for OSQ.

\paragraph{Cost per Sample}
\begin{equation}
\text{Cost}_{\text{sample}}^{(m)} = \frac{\text{Tokens}_{\text{avg}}^{(m)} \times \text{price}_{\text{output}}^{(m)}}{10^6}
\end{equation}
where $\text{price}_{\text{output}}^{(m)}$ is output token price in \$/million tokens.

\paragraph{Accuracy per Dollar}
\begin{equation}
\text{Efficiency}^{(m)} = \frac{\text{Acc}^{(m)}}{\text{Cost}_{\text{sample}}^{(m)}}
\end{equation}

Higher values indicate more cost-efficient models.

\paragraph{ROI for Reasoning Models}
Compares thinking model $m_t$ to baseline model $m_b$:
\begin{equation}
\text{ROI} = \frac{\Delta \text{Acc}}{\Delta \text{Cost}} = \frac{\text{Acc}^{(m_t)} - \text{Acc}^{(m_b)}}{\text{Cost}_{\text{sample}}^{(m_t)} - \text{Cost}_{\text{sample}}^{(m_b)}}
\end{equation}

Interpretation: ROI represents additional accuracy gained per additional dollar spent.

\paragraph{Break-Even Analysis}
At what benchmark scale $N$ does the total cost of thinking model equal that of baseline model for equivalent accuracy gain?

\begin{equation}
N_{\text{breakeven}} = \frac{\text{Fixed overhead cost}}{\text{Cost}_{\text{sample}}^{(m_t)} - \text{Cost}_{\text{sample}}^{(m_b)}}
\end{equation}

\subsubsection{Thinking Model Token Accounting}

Reasoning-extended models (e.g., DeepSeek-R1, QwQ-32B) produce outputs with explicit reasoning enclosed in special tags:
\begin{verbatim}
<think>
[Extended reasoning process...]
</think>

[Final answer]
\end{verbatim}

Token accounting separates:
\begin{itemize}
    \item \textbf{Reasoning tokens}: Content within \texttt{<think>} tags
    \item \textbf{Answer tokens}: Content outside \texttt{<think>} tags
\end{itemize}

This enables analysis of whether reasoning token expenditure correlates with accuracy improvement.

\subsubsection{Visualizations}

\begin{itemize}
    \item \textbf{Cost vs. Accuracy scatter}: Models plotted by cost/sample (x-axis) and accuracy (y-axis)
    \item \textbf{Efficiency frontier}: Convex hull of models representing Pareto-optimal cost-accuracy tradeoffs
    \item \textbf{ROI curves}: ROI as function of benchmark scale for different model pairs
    \item \textbf{Token distribution histograms}: Distribution of tokens/sample for thinking vs. standard models
    \item \textbf{Break-even plots}: Break-even point visualized as intersection of cost curves
\end{itemize}

Example cost-accuracy scatter in Figure~\ref{fig:cost_accuracy}.

\begin{figure}[htbp]
    \centering
    % TODO: Include actual Phase 6 cost-accuracy scatter
    \includegraphics[width=0.8\textwidth]{figures/cost_accuracy.pdf}
    \caption{Cost-accuracy scatter plot with efficiency frontier. Each point represents a model, positioned by average cost per sample (x-axis) and accuracy (y-axis). The efficiency frontier (solid line) connects Pareto-optimal models. Models on the frontier represent best cost-accuracy tradeoffs; models below are dominated by frontier models.}
    \label{fig:cost_accuracy}
\end{figure}

\subsubsection{Example Cost Analysis}

Table~\ref{tab:cost_example} presents hypothetical cost analysis for comparison (actual data in results chapter).

\begin{table}[htbp]
    \centering
    \caption{Example tokenomics comparison between standard and reasoning-extended models.}
    \label{tab:cost_example}
    \begin{tabular}{lrrrrr}
        \toprule
        \textbf{Model} & \textbf{Acc} & \textbf{Avg Tokens} & \textbf{Cost/Sample} & \textbf{Acc/\$} & \textbf{ROI} \\
        \midrule
        Llama 3.2 3B & 0.61 & 187 & \$0.00002 & 30,500 & -- \\
        DeepSeek-R1 7B & 0.78 & 3,245 & \$0.00097 & 803 & 179 \\
        GPT-4o & 0.89 & 412 & \$0.02060 & 43 & 13.6 \\
        \bottomrule
    \end{tabular}
\end{table}

Interpretation: DeepSeek-R1 achieves 17\% higher accuracy than Llama 3.2 at substantially higher cost (\$0.00095/sample more), yielding ROI = 179 (gain 179 percentage points accuracy per additional dollar). GPT-4o provides best absolute accuracy but lowest efficiency (43 accuracy per dollar) and lowest ROI relative to baseline.

%============================================================================
\section{Summary}

This chapter detailed the six-phase methodology for comprehensive evaluation of large language models on systems engineering knowledge:

\begin{enumerate}
    \item \textbf{Phase 1}: Acquired SysEngBench dataset (1,144 MCQs) and characterized distribution across INCOSE knowledge areas
    \item \textbf{Phase 2}: Developed two-stage LLM-based pipeline to convert 845 suitable MCQs to OSQ format with grading rubrics
    \item \textbf{Phase 3}: Generated four position-rotated variants for position bias detection
    \item \textbf{Phase 4}: Executed model inference across six evaluation tasks (base MCQ, 4 position variants, OSQ) using lm-evaluation-harness
    \item \textbf{Phase 5}: Applied four LLM-as-a-judge approaches to score OSQ responses
    \item \textbf{Phase 6}: Performed statistical analysis covering position bias, format comparison, and tokenomics with publication-ready visualizations
\end{enumerate}

All datasets, code, and evaluation outputs are publicly available for reproducibility. The following chapters present results (Chapter~\ref{ch:results}), discuss implications (Chapter~\ref{ch:discussion}), and conclude with recommendations for practice and future research (Chapter~\ref{ch:conclusion}).

%============================================================================
% Appendices would be included here if this were a standalone document,
% but typically dissertations place appendices at the end after all chapters

% Example placeholder for appendices references:
% \appendix
% \section{Conversion and Judging Prompts}\label{app:conversion_prompts}
% \section{Complete Model List}\label{app:model_list}
% \section{Dataset Visualizations}\label{app:dataset_viz}
% \section{Statistical Test Results Tables}\label{app:stat_tables}
